\section*{致谢}
我们感谢徐峰、袁立、Shi-Cheng Xia、Jianhui Zhang 与 Yong Zhao 的有益讨论。
我们感谢 CLS 合作组分享用于本研究的格点系综。
%
格点 QCD 计算使用 Chroma 软件套件完成~\cite{Edwards:2004sx}。
%
数值计算得到中国科学院战略性先导科技专项（Grant No.\ XDC01040100）、ITP-CAS 高性能计算集群以及江苏大规模特殊计算平台的支持。数值模拟的部分设置在上海交通大学高性能计算中心支持的 $\pi$ 2.0 集群上进行。
J.~Hua 受 NSFC Grant No.\ 11735010 与 11947215 资助。
%
Y.-S.~Liu 受国家自然科学基金 Grant No.\ 11905126 资助。
%
M.~Schlemmer 与 A.~Sch\"afer 受 DFG 合作研究中心 CRC/TRR-55 支持。
%
P.~Sun 受国家自然科学基金 Grant No.\ 11975127 及江苏省特聘教授计划资助。
%
W.~Wang 部分受国家自然科学基金 Grant No.\ 11735010、11911530088 以及上海市自然科学基金 Grant No.\ 15DZ2272100 资助。
%
Q.-A.~Zhang 受中国博士后科学基金会与博士后创新人才支持计划（Grant No.\ BX20190207）资助。

\section{补充材料}


\subsection{模拟检验}\label{subsec:simchecks}

\begin{figure}[!th]
\begin{center}
\includegraphics[width=0.5\textwidth]{images/06_eff_mass.pdf}
\caption{ 由二点函数π介子质量给出的π介子态色散关系。直到 $8\pi/L$（$\sim$2 GeV）的数据可用公式 $E_{\pi}=\sqrt{m_{\pi}^2+c_1 P^2+c_2P^4a^2}$ 描述，其中 $c_1=0.9945(40)$、$c_2=-0.0282(27)$。$8\pi/L$ 处相对连续极限的偏离约为 $2\%$。
}\label{fig:dispersion}
\end{center}
\end{figure}

图~\ref{fig:dispersion} 展示了我们所用π介子质量下的色散关系。曲线为基于公式 $E_{\pi}=\sqrt{m_{\pi}^2+c_1 P^2+c_2P^4a^2}$ 的拟合，根号内最后一项参数化离散误差。我们使用的动量最高达 $8\pi/L$（$\sim$2 GeV）。拟合给出 $c_1=0.9945(40)$、$c_2=-0.0282(27)$，与由二点函数计算的基态能量一致。这表明离散误差很小。因此可以预期，在连续极限下色散关系将恢复标准的 $E_{\pi}=\sqrt{m_{\pi}^2+P^2}$。

\begin{figure}[!th]
\begin{center}
\includegraphics[width=0.4\textwidth]{images/07a_hist_ori3pt_P3_b3_tsep8.pdf}
\includegraphics[width=0.4\textwidth]{images/07b_bin_3pt_P3_b3_tsep8.pdf}
\caption{模拟的统计检验，以 $P^z=6\pi/L$、$b_{\perp}$=3、$t_{2}$=8、$t=t_{2}/2$ 的形状因子为例。左图为{纳入后续分析的 $864$ 个组态}的直方图，右图为对所有组态上的测量取平均后的 bin 尺寸依赖。}\label{fig:statistics}
\end{center}
\end{figure}



以 $P^z=6\pi/L$、$b_{\perp}$=3、$t_{2}$=8、$t=t_{2}/2$ 的形状因子为例，图~\ref{fig:statistics} 展示了我们所做测量的统计检验。{我们分析了 868 个组态，其中 4 个因存在很强的局域伪迹而被剔除。}左图为 864（组态）$\times$ 48（时间{片}）次测量的直方图。{已经注意到，在动力学组态上使用轻质量 clover 作用量和（或）较粗的格点时，尽管非常罕见，仍可能出现反常测量，因为临界点并不十分稳定。事实证明，某些很强的局域伪迹并未在更细格距的其他 CLS 系综中观察到，但可能出现在粗格距系综的少数组态中，例如本文分析所用的 A654 系综。由于这只占总组态的一小部分（即 $4/868$），剔除这些组态是合理的。}

 对同一组态上的测量取平均后，我们发现自关联效应可以忽略，因为结果没有明显的 bin 尺寸依赖，如图~\ref{fig:statistics} 右图所示。

\subsection{TMDWF 的 $\ell$ 依赖}\label{subsec:ldep}


\begin{figure}[!th]
\begin{center}
\includegraphics[width=0.9\textwidth]{images/08_tmdwf_4fig_mod.pdf}
\caption{$|\phi_\ell(0,b_\perp,P^z,\ell)|$ 的 $\ell$ 依赖图：$\{P^z, b_{\perp}, t\}=\{6\pi/L, 3a, 6a\}$（左上，即图~\ref{fig:L_dependence} 所示情形）、$\{P^z, b_{\perp}, t\}=\{8\pi/L, 3a, 6a\}$（右上）、$\{P^z, b_{\perp}, t\}=\{6\pi/L, 6a, 6a\}$（左下）、$\{P^z, b_{\perp}, t\}=\{6\pi/L, 3a, 8a\}$（右下）。}\label{fig:L_dependence2}
\end{center}
\end{figure}

%%%%%%%%%%


在图~\ref{fig:L_dependence2} 中，我们给出了另外几个情形下 $|\phi_\ell(0,b_\perp,P^z,\ell)|$ 的 $\ell$ 依赖，它们与图~\ref{fig:L_dependence} 所示 $\{P^z, b_{\perp}, t\}=\{6\pi/L, 3a, 6a\}$ 情形类似，但 $P^z$、$b_{\perp}$ 与 $t$ 更大。


\begin{figure}[!th]
\begin{center}
\includegraphics[width=0.45\textwidth]{images/09a_HeatMap_b0.pdf}
\includegraphics[width=0.45\textwidth]{images/09b_HeatMap_b3.pdf}
\includegraphics[width=0.45\textwidth]{images/09c_HeatMap_b6.pdf}
\includegraphics[width=0.45\textwidth]{images/09d_heatmap_of_renormTMDWF.pdf}
\caption{{$P=1.58$ GeV 时非微扰重正化因子的热图（$b_{\perp}=0$：左上，$b_{\perp}=3a$：右上，$b_{\perp}=6a$：左下）与矩阵元（右下，按 $b_{\perp}=\{0,~3a,~6a\}$ 的顺序展示）。图中 $\Gamma$ 表示算符中的洛伦兹结构，${\cal P}$ 表示投影。}}\label{fig:HeatMap_NPR}
\end{center}
\end{figure}


\subsection{TMDWF 算符混合的估计}\label{subsec:mixing}

{
为了估计算符混合效应，我们采用与文献~\cite{Shanahan:2019zcq,Shanahan:2020zxr}相同的方法，计算非微扰 RI/MOM 重正化/混合因子，
\begin{align}
{\cal M}_{{\cal P}\Gamma}=\textrm{Tr}\left[{\cal P}\Big\langle q(p)\Big|\overline q_1\left(\frac{z }{2}n^z +\vec b\right)\Gamma \,{\cal W}(\vec b,\ell)q_2\left(-\frac{z}{2}n^z \right)\Big| q(p)\Big\rangle\right],
\end{align}
其中 $\Gamma$ 表示算符中的洛伦兹结构，${\cal P}$ 表示投影。相对混合效应用比值 ${\cal M}_{{\cal P}\Gamma}/{\cal M}_{\Gamma\Gamma}$ 来衡量。三个横向间隔 $b_\perp=(0a, 3a, 6a)\hat{n}_x$ 与离壳夸克动量 $p=(p_t,p_x,p_y,p_z)=(3.16,0,1.58,0)$ GeV 的结果以热图形式示于图~\ref{fig:HeatMap_NPR}（分别对应左上、右上与左下小图）。混合效应随 $b_\perp$ 增大而增强。这一模式与文献~\cite{Shanahan:2019zcq} 的微扰计算一致。}

{
准 TMDWF 中的相对混合效应可以通过具有给定洛伦兹结构 ${\cal P}$ 的裸准 TMDWF 与相应混合因子 ${\cal M}_{\Gamma{\cal P}}$ 的乘积来估计，
\begin{align}
\frac{\delta_{{\cal P}}\phi(z,b_\perp,P^z)}{\phi(z,b_\perp,P^z)}\simeq\frac{{\cal M}_{\Gamma_{\phi}{\cal P}}\Big\langle 0\Big|\overline q_1\left(\frac{z }{2}n^z +\vec b\right){\cal P}\,{\cal W}(\vec b,\ell)q_2\left(-\frac{z}{2}n^z \right)\Big| {\pi(\vec{P})}\Big\rangle}{\textrm{Re}[{\cal M}_{\Gamma_{\phi}\Gamma_{\phi}}\Big\langle 0\Big|\overline q_1\left(\frac{z }{2}n^z +\vec b\right){\Gamma_{\phi}}\,{\cal W}(\vec b,\ell)q_2\left(-\frac{z}{2}n^z \right)\Big| {\pi(\vec{P})}\Big\rangle]},
\end{align}
其中 $\Gamma_{\phi}=\gamma_5\gamma_t$。
对于 $b_\perp=(0a, 3a, 6a)\hat{n}_x$，结果在图~\ref{fig:HeatMap_NPR} 右下小图中给出。从图可见，对横向间隔 $b\sim$ 0.6 fm，算符混合效应可达 $5\%$ 量级，而对更小的横向间隔则不那么显著。
}


%%%%%%%%%%%%%%%%%%


\subsection{内禀软函数的列表结果}\label{subsec:tabulated}

\begin{figure}[!th]
\includegraphics[width=0.45\textwidth]{images/04_sf_perturb_sqrt2.pdf}
\includegraphics[width=0.45\textwidth]{images/14_sf_normalscale.pdf}
\caption{ 按式{$S_{I,\overline{\rm MS}}(b_\perp,\mu)=\frac{F(b_\perp,P^z)}{F(b_{\perp,0},P^z)}\frac{|\phi(0,b_{\perp,0},P^z)|^2}{|\phi(0,b_\perp,P^z)|^2}+{\cal O}(\alpha_s, \gamma^{-2})$} 取 $b_{\perp,0}=a$ 时内禀软因子随 $b_{\perp}$ 的变化。不同π介子动量 $P^z$ 的结果彼此一致。虚线为一圈计算结果，即 {$S_{I,\overline{\rm MS}}(b_{{\perp}},\mu)=1-\frac{\alpha_sC_F}{\pi}\ln\frac{\mu^2 b_{{\perp}}^2}{4 e^{-2\gamma_E}}+{\cal O}(\alpha_s)$}，强耦合常数取 $\alpha_s(1/b_{\perp})$。{阴影带对应 $\alpha_s$ 的标度不确定度：$\mu\in [1/\sqrt{2},\sqrt{2}]\times 1/b_\perp$。来自算符混合的系统不确定度已被计入。{两幅小图显示相同的数据，左图为对数标度，右图为普通标度。}}}\label{fig:quasi-soft-factor2}
\end{figure}

{如图~\ref{fig:quasi-soft-factor2} 所示，为便于比较，我们在表~\ref{table:soft-function} 中列表给出了内禀软函数的结果。考虑到强耦合常数 $\alpha_s(\mu)$ 的标度{依赖}带来的误差（$\mu\in [1/\sqrt{2},\sqrt{2}]\times 1/b_\perp$），$b\le3a$ 的微扰结果与我们的计算一致。}


\begin{table}[htbp]
\begin{center}
\caption{图~\ref{fig:quasi-soft-factor} 所示内禀软函数的列表结果。}\label{table:soft-function}
\begin{tabular}{|c|ccccccc|cc}
\hline
 & $b=a$         & $b=2a$         & $b=3a$         & $b=4a$         & $b=5a$         & $b=6a$         & $b=7a$         \\
\hline
$P=1.05$ GeV                   & 1.000(8)  & 0.567(7)  & 0.343(6)  & 0.224(6)  & 0.153(6)  & 0.106(7)  & 0.071(7)  \\
$P=1.58$ GeV     & 1.000(20) & 0.557(17) & 0.329(13) & 0.209(14) & 0.142(18) & 0.099(22) & 0.063(25) \\
$P=2.11$ GeV & 1.000(29) & 0.571(62) & 0.374(63) & 0.223(52) & 0.119(50) & 0.043(46) & 0.047(84)  \\
\hline
%pQCD [1/2, 2] & $1.000_{-0.009}^{+0.020}$   &  $0.703_{-0.326}^{+0.093}$  &  $0.303_{-3.032}^{+0.298}$  & -  & - & - & - \\
pQCD & $1.000_{-0.008}^{+0.005}$   &  $0.703_{-0.100}^{+0.057}$  &  $0.303_{-0.461}^{+0.193}$  & $-0.334_{-2.142}^{+0.498}$  & - & - & - \\
\hline
\end{tabular}
\end{center}
\end{table}

\subsection{形状因子的两态拟合}\label{subsec:twostate}


本工作中，我们执行如下联合拟合以获得减除准 TMDWF 的模 $|\phi_\ell(0,b_\perp,P^z,\ell)|$ 与软因子 $S_I(b_\perp)$（取 $\ell=7a$），
\begin{align}
&\frac{C_3(b_\perp,P^z,t_{\rm sep},t)}{C_2(0,P^z,0,t_{\rm sep})}=\frac{|\tilde{\phi}_\ell(0,b_\perp,P^z,\ell)|^2\tilde{S}_I(b_\perp)+C_1(e^{-\Delta E t}+e^{-\Delta E (t_{\rm sep}-t)})+C_2e^{-\Delta E t_{\rm sep}}}{1+C_0 e^{-\Delta E t_{\rm sep}}}, \nonumber\\
&\frac{C_2(b_\perp,P^z,0,t)}{C_2(0,P^z,0,t)}=\frac{|\tilde{\phi}_\ell(0,b_\perp,P^z,\ell)|e^{\theta(b_\perp,P^z,\ell)}(1+C_3e^{-\Delta E t})}{1+C_0e^{-\Delta E t}},
\label{eq:fit}
\end{align}
其中
\begin{align}
\tilde{\phi}_\ell(0,b_\perp,P^z,\ell)=\frac{\phi_\ell(0,b_\perp,P^z,\ell)}{\phi_\ell(0,0,P^z,0)},\ \tilde{S}_I(b_\perp)=\frac{{\phi}_L(0,0,P^z,0)A_w}{2EA_p}S_I(b_\perp),
\end{align}
而 $\theta(b_\perp,P^z,\ell)$ 是准 TMDWF 的相位。
%
$\tilde{S}_I$ 定义中的附加因子将在 $\tilde{S}_I(b_\perp=a)$ 中被抵消{，当我们考虑如下比值时，
\begin{align}\label{eq:S_ratio_app}
S_{I,\overline{\rm MS}}(b_\perp,\mu)= \left(\frac{S_I(b_{\perp},1/a)}{S_I(b_{\perp,0},1/a)}\right)S_{I,\overline{\rm MS}}(b_{\perp,0},\mu).
\end{align}}

%%%%%%%%%%%%%%%%
\begin{figure}[!th]
\begin{center}
\includegraphics[width=0.45\textwidth]{images/10a_ff_P3b1.pdf}
\includegraphics[width=0.45\textwidth]{images/03_ff_P3b3.pdf}\\
\includegraphics[width=0.45\textwidth]{images/10c_ff_P3b5.pdf}
\includegraphics[width=0.45\textwidth]{images/10d_test_ff_P3b1.pdf}
\caption{ 比值 $C_3(b_\perp,P^z,t_{\rm sep},t)/C_2(0, P^z,0,t_{\rm sep})$ 作为 $t_{\rm sep}$ 与 $t$ 的函数：$\{P^z, b_{\perp}\}=\{6\pi/L, 1a\}$（左上）、$\{P^z, b_{\perp}\}=\{6\pi/L, 3a\}$（右上）与 $\{P^z, b_{\perp}\}=\{6\pi/L, 5a\}$（左下）。{右下图中，我们对 $\{P^z, b_{\perp}\}=\{6\pi/L, 1a\}$ 情形剔除了 $t_{\rm sep}=6$a 的数据，发现拟合结果{与左上情形一致}。}}\label{fig:joint_fit2}
\end{center}
\end{figure}
%%%%%%%%%%%%%%%%

在图~\ref{fig:joint_fit2} 中，我们{展示了} $P^z=6\pi/L$、$b_{\perp}=\{1a,3a,5a\}$ 的比值 $C_3(b_\perp,P^z,t_{\rm sep},t)/C_2(0, P^z,0,t_{\rm sep})$，并与两态拟合预言（彩色带）及拟合的基态贡献（灰带）比较。所有情形的数据与拟合都符合得很好。
{这一致性表明来自拟合窗口的系统不确定度是温和的。作为另一项估计，我们对 $\{P^z, b_{\perp}\}=\{6\pi/L, 1a\}$ 情形剔除了 $t_{\rm sep}=6$a 的数据，结果示于图~\ref{fig:joint_fit2} 右下图。可以发现拟合结果{在不确定度范围内与左上图情形一致。}}


%%%%%%%%%%%%%%%%
\begin{figure}[!th]
\begin{center}
\includegraphics[width=0.7\textwidth]{images/11_differential_summed_ratio.pdf}
\caption{ $\{P^z, b_{\perp}\}=\{6\pi/L, 3a\}$ 时比值 $C_3(b_\perp,P^z,t_{\rm sep},t)/C_2(0, P^z,0,t_{\rm sep})$ 随 $t_{\rm sep}$ 与 $t$ 的变化，并与差分求和比 $R(b_\perp,P^z,t_{\rm sep})$ 比较。}\label{fig:joint_fit3}
\end{center}
\end{figure}
%%%%%%%%%%%%%%%%

作为又一项检验，我们还考虑差分求和比
\begin{align}
R(b_\perp,P^z,t_{\rm sep})&\equiv SR(b_\perp,P^z,t_{\rm sep})-SR(b_\perp,P^z,t_{\rm sep}-1)=|\tilde{\phi}_\ell(0,b_\perp,P^z,\ell)|^2\tilde{S}_I(b_\perp)+{\cal O}(e^{-\Delta E t_{\rm sep}}),\ \nonumber\\
SR(b_\perp,P^z,t_{\rm sep})&\equiv \sum_{0<t<t_{\rm sep}} \frac{C_3(b_\perp,P^z,t_{\rm sep},t)}{C_2(0, P^z,0,t_{\rm sep})}.
\end{align}
作为一个例子，图~\ref{fig:joint_fit3} 中对 $\{P^z, b_{\perp}\}=\{6\pi/L, 3a\}$ 绘出了 $R(b_\perp,P^z,t_{\rm sep})$ 随 $t_{\rm sep}$ 的变化，并与标准两态拟合比较。可以看到，在大 $t_{\rm sep}$ 处 $R(b_\perp,P^z,t_{\rm sep})$ 与两态拟合给出的基态贡献相符。


\subsection{提取 Collins-Soper 核时可能出现的虚部}\label{subsec:imagpart}

\begin{figure}[!th]
\begin{center}
\includegraphics[width=0.44\textwidth]{images/12a_CS_kernel_re_for_compare.pdf}
\includegraphics[width=0.44\textwidth]{images/12b_CS_kernel_im.pdf}
\caption{ 采用近似 $C(xP^z,\mu)=1+{\cal O}(\alpha_s)$ 并按式~(\ref{eq:CSK_full}) 定义时，Collins-Soper 核的实部（左）与虚部（右）。}\label{fig:imaginal}
\end{center}
\end{figure}

%%%%%%%%%%%%%
\begin{figure}[!th]
\begin{center}
\includegraphics[width=0.45\textwidth]{images/13_CS_kernel_compare.pdf}
\caption{$K$ 与 $-|K'|=-\sqrt{K^2+K'^2_{\rm Im}}$ 的比较。二者彼此一致。}\label{fig:compare}
\end{center}
\end{figure}

按如下定义的 Collins-Soper 核
\begin{align}
&K'(b_\perp,\mu)=\frac{1}{\ln(P_1^z/P_2^z)}\ln\frac{C(xP_2^z,\mu) \Phi_{\overline{\rm MS}}(x,b_\perp,P_1^z,\mu)}{C(xP_1^z,\mu) \Phi_{\overline{\rm MS}}(x,b_\perp,P_2^z,\mu)}+{\cal O}(\gamma^{-2})=\frac{1}{\ln(P_1^z/P_2^z)}\ln\frac{\phi(0,b_\perp,P_1^z)}{\phi(0,b_\perp,P_2^z)}+{\cal O}(\alpha_s, \gamma^{-2})\, \nonumber\\
&\quad\quad\quad\ \ =\frac{1}{\ln(P_1^z/P_2^z)}\ln\left|\frac{\phi(0,b_\perp,P_1^z)}{\phi(0,b_\perp,P_2^z)}\right|+\frac{i}{\ln(P_1^z/P_2^z)}\left(\theta(0,b_\perp,P_1^z)-\theta(0,b_\perp,P_2^z)\right)+{\cal O}(\alpha_s, \gamma^{-2})\, \label{eq:CSK_full}
\end{align}
应当是实的，但当采用近似 $C(xP^z,\mu)=1+{\cal O}(\alpha_s)$ 时，相位 $\theta(0,b_\perp,P^z)=\textrm{tan}^{-1}\frac{\phi_{\rm Im}(0,b_\perp,P^z)}{\phi_{\rm Re}(0,b_\perp,P^z)}$ 的 $P^z$ 依赖会给 $K'$ 引入一个{虚部}。图~\ref{fig:imaginal} 展示了两种 $P_1/P_2$ 组合下实部与虚部随 $b_\perp$ 的变化。实部 $K'_{\rm Re}=K$（左图）即正文所用的定义。非零的虚部（右图）反映了匹配不精确带来的系统不确定度。如图~\ref{fig:compare} 所示，在 $K$ 的统计不确定度范围内，$K$ 与 $-|K'|=-\sqrt{K^2+K'^2_{\rm Im}}$ 仍然是一致的。

为估计匹配不精确的影响，我们把 $|K'_{\rm Im}|$ 视为一个系统不确定度，并与 $K$ 的统计不确定度平方相加。


\begin{thebibliography}{99}
%\cite{Ellis:1991qj}
\bibitem{Ellis:1991qj} 
  R.~K.~Ellis, W.~J.~Stirling and B.~R.~Webber,
  %``QCD and collider physics,''
  Camb.\ Monogr.\ Part.\ Phys.\ Nucl.\ Phys.\ Cosmol.\  {\bf 8}, 1 (1996).
  %%CITATION = CMPCE,8,1;%%
  %730 citations counted in INSPIRE as of 08 Oct 2020


%\cite{Lin:2017snn}
\bibitem{Lin:2017snn} 
  H.~W.~Lin {\it et al.},
  %``Parton distributions and lattice QCD calculations: a community white paper,''
  Prog.\ Part.\ Nucl.\ Phys.\  {\bf 100}, 107 (2018)
  doi:10.1016/j.ppnp.2018.01.007
  [arXiv:1711.07916 [hep-ph]].
  %%CITATION = doi:10.1016/j.ppnp.2018.01.007;%%
  %125 citations counted in INSPIRE as of 08 Oct 2020


%\cite{Collins:1981uk}
\bibitem{Collins:1981uk} 
  J.~C.~Collins and D.~E.~Soper,
  %``Back-To-Back Jets in QCD,''
  Nucl.\ Phys.\ B {\bf 193}, 381 (1981)
  Erratum: [Nucl.\ Phys.\ B {\bf 213}, 545 (1983)].
  doi:10.1016/0550-3213(81)90339-4
  %%CITATION = doi:10.1016/0550-3213(81)90339-4;%%
  %1218 citations counted in INSPIRE as of 08 Oct 2020


%\cite{Collins:1984kg}
\bibitem{Collins:1984kg} 
  J.~C.~Collins, D.~E.~Soper and G.~F.~Sterman,
  %``Transverse Momentum Distribution in Drell-Yan Pair and W and Z Boson Production,''
  Nucl.\ Phys.\ B {\bf 250}, 199 (1985).
  doi:10.1016/0550-3213(85)90479-1
  %%CITATION = doi:10.1016/0550-3213(85)90479-1;%%
  %1252 citations counted in INSPIRE as of 08 Oct 2020


%\cite{Ji:2004wu}
\bibitem{Ji:2004wu} 
  X.~d.~Ji, J.~p.~Ma and F.~Yuan,
  %``QCD factorization for semi-inclusive deep-inelastic scattering at low transverse momentum,''
  Phys.\ Rev.\ D {\bf 71}, 034005 (2005)
  doi:10.1103/PhysRevD.71.034005
  [hep-ph/0404183].
  %%CITATION = doi:10.1103/PhysRevD.71.034005;%%
  %601 citations counted in INSPIRE as of 08 Oct 2020


%\cite{Ji:2004xq}
\bibitem{Ji:2004xq} 
  X.~d.~Ji, J.~P.~Ma and F.~Yuan,
  %``QCD factorization for spin-dependent cross sections in DIS and Drell-Yan processes at low transverse momentum,''
  Phys.\ Lett.\ B {\bf 597}, 299 (2004)
  doi:10.1016/j.physletb.2004.07.026
  [hep-ph/0405085].
  %%CITATION = doi:10.1016/j.physletb.2004.07.026;%%
  %359 citations counted in INSPIRE as of 08 Oct 2020


%\cite{Echevarria:2015byo}
\bibitem{Echevarria:2015byo} 
  M.~G.~Echevarria, I.~Scimemi and A.~Vladimirov,
  %``Universal transverse momentum dependent soft function at NNLO,''
  Phys.\ Rev.\ D {\bf 93}, no. 5, 054004 (2016)
  doi:10.1103/PhysRevD.93.054004
  [arXiv:1511.05590 [hep-ph]].
  %%CITATION = doi:10.1103/PhysRevD.93.054004;%%
  %76 citations counted in INSPIRE as of 08 Oct 2020


%\cite{Li:2016ctv}
\bibitem{Li:2016ctv} 
  Y.~Li and H.~X.~Zhu,
  %``Bootstrapping Rapidity Anomalous Dimensions for Transverse-Momentum Resummation,''
  Phys.\ Rev.\ Lett.\  {\bf 118}, no. 2, 022004 (2017)
  doi:10.1103/PhysRevLett.118.022004
  [arXiv:1604.01404 [hep-ph]].
  %%CITATION = doi:10.1103/PhysRevLett.118.022004;%%
  %100 citations counted in INSPIRE as of 08 Oct 2020


%\cite{Ji:2019sxk}
\bibitem{Ji:2019sxk} 
  X.~Ji, Y.~Liu and Y.~S.~Liu,
  %``TMD soft function from large-momentum effective theory,''
  Nucl.\ Phys.\ B {\bf 955}, 115054 (2020)
  doi:10.1016/j.nuclphysb.2020.115054
  [arXiv:1910.11415 [hep-ph]].
  %%CITATION = doi:10.1016/j.nuclphysb.2020.115054;%%
  %19 citations counted in INSPIRE as of 08 Oct 2020


%\cite{Ji:2013dva}
\bibitem{Ji:2013dva} 
  X.~Ji,
  %``Parton Physics on a Euclidean Lattice,''
  Phys.\ Rev.\ Lett.\  {\bf 110}, 262002 (2013)
  doi:10.1103/PhysRevLett.110.262002
  [arXiv:1305.1539 [hep-ph]].
  %%CITATION = doi:10.1103/PhysRevLett.110.262002;%%
  %368 citations counted in INSPIRE as of 08 Oct 2020


%\cite{Ji:2014gla}
\bibitem{Ji:2014gla} 
  X.~Ji,
  %``Parton Physics from Large-Momentum Effective Field Theory,''
  Sci.\ China Phys.\ Mech.\ Astron.\  {\bf 57}, 1407 (2014)
  doi:10.1007/s11433-014-5492-3
  [arXiv:1404.6680 [hep-ph]].
  %%CITATION = doi:10.1007/s11433-014-5492-3;%%
  %157 citations counted in INSPIRE as of 08 Oct 2020


%\cite{Collins:2011zzd}
\bibitem{Collins:2011zzd} 
  J.~Collins,
  %``Foundations of perturbative QCD,''
  Camb.\ Monogr.\ Part.\ Phys.\ Nucl.\ Phys.\ Cosmol.\  {\bf 32}, 1 (2011).
  %%CITATION = CMPCE,32,1;%%
  %450 citations counted in INSPIRE as of 08 Oct 2020


%\cite{Bauer:2000yr}
\bibitem{Bauer:2000yr} 
  C.~W.~Bauer, S.~Fleming, D.~Pirjol and I.~W.~Stewart,
  %``An Effective field theory for collinear and soft gluons: Heavy to light decays,''
  Phys.\ Rev.\ D {\bf 63}, 114020 (2001)
  doi:10.1103/PhysRevD.63.114020
  [hep-ph/0011336].
  %%CITATION = doi:10.1103/PhysRevD.63.114020;%%
  %1421 citations counted in INSPIRE as of 08 Oct 2020


%\cite{Bauer:2001ct}
\bibitem{Bauer:2001ct} 
  C.~W.~Bauer and I.~W.~Stewart,
  %``Invariant operators in collinear effective theory,''
  Phys.\ Lett.\ B {\bf 516}, 134 (2001)
  doi:10.1016/S0370-2693(01)00902-9
  [hep-ph/0107001].
  %%CITATION = doi:10.1016/S0370-2693(01)00902-9;%%
  %735 citations counted in INSPIRE as of 08 Oct 2020


%\cite{Bauer:2001yt}
\bibitem{Bauer:2001yt} 
  C.~W.~Bauer, D.~Pirjol and I.~W.~Stewart,
  %``Soft collinear factorization in effective field theory,''
  Phys.\ Rev.\ D {\bf 65}, 054022 (2002)
  doi:10.1103/PhysRevD.65.054022
  [hep-ph/0109045].
  %%CITATION = doi:10.1103/PhysRevD.65.054022;%%
  %1228 citations counted in INSPIRE as of 08 Oct 2020


%\cite{Ji:2020ect}
\bibitem{Ji:2020ect} 
  X.~Ji, Y.~S.~Liu, Y.~Liu, J.~H.~Zhang and Y.~Zhao,
  %``Large-Momentum Effective Theory,''
  arXiv:2004.03543 [hep-ph].
  %%CITATION = ARXIV:2004.03543;%%
  %32 citations counted in INSPIRE as of 08 Oct 2020


%\cite{Cichy:2018mum}
\bibitem{Cichy:2018mum} 
  K.~Cichy and M.~Constantinou,
  %``A guide to light-cone PDFs from Lattice QCD: an overview of approaches, techniques and results,''
  Adv.\ High Energy Phys.\  {\bf 2019}, 3036904 (2019)
  doi:10.1155/2019/3036904
  [arXiv:1811.07248 [hep-lat]].
  %%CITATION = doi:10.1155/2019/3036904;%%
  %66 citations counted in INSPIRE as of 08 Oct 2020


%\cite{Ji:2014hxa}
\bibitem{Ji:2014hxa} 
  X.~Ji, P.~Sun, X.~Xiong and F.~Yuan,
  %``Soft factor subtraction and transverse momentum dependent parton distributions on the lattice,''
  Phys.\ Rev.\ D {\bf 91}, 074009 (2015)
  doi:10.1103/PhysRevD.91.074009
  [arXiv:1405.7640 [hep-ph]].
  %%CITATION = doi:10.1103/PhysRevD.91.074009;%%
  %67 citations counted in INSPIRE as of 08 Oct 2020


%\cite{Ji:2018hvs}
\bibitem{Ji:2018hvs} 
  X.~Ji, L.~C.~Jin, F.~Yuan, J.~H.~Zhang and Y.~Zhao,
  %``Transverse momentum dependent parton quasidistributions,''
  Phys.\ Rev.\ D {\bf 99}, no. 11, 114006 (2019)
  doi:10.1103/PhysRevD.99.114006
  [arXiv:1801.05930 [hep-ph]].
  %%CITATION = doi:10.1103/PhysRevD.99.114006;%%
  %41 citations counted in INSPIRE as of 08 Oct 2020


%\cite{Ebert:2018gzl}
\bibitem{Ebert:2018gzl} 
  M.~A.~Ebert, I.~W.~Stewart and Y.~Zhao,
  %``Determining the Nonperturbative Collins-Soper Kernel From Lattice QCD,''
  Phys.\ Rev.\ D {\bf 99}, no. 3, 034505 (2019)
  doi:10.1103/PhysRevD.99.034505
  [arXiv:1811.00026 [hep-ph]].
  %%CITATION = doi:10.1103/PhysRevD.99.034505;%%
  %40 citations counted in INSPIRE as of 08 Oct 2020


%\cite{Ebert:2019okf}
\bibitem{Ebert:2019okf} 
  M.~A.~Ebert, I.~W.~Stewart and Y.~Zhao,
  %``Towards Quasi-Transverse Momentum Dependent PDFs Computable on the Lattice,''
  JHEP {\bf 1909}, 037 (2019)
  doi:10.1007/JHEP09(2019)037
  [arXiv:1901.03685 [hep-ph]].
  %%CITATION = doi:10.1007/JHEP09(2019)037;%%
  %33 citations counted in INSPIRE as of 08 Oct 2020


%\cite{Ji:2019ewn}
\bibitem{Ji:2019ewn} 
  X.~Ji, Y.~Liu and Y.~S.~Liu,
  %``Transverse-Momentum-Dependent PDFs from Large-Momentum Effective Theory,''
  arXiv:1911.03840 [hep-ph].
  %%CITATION = ARXIV:1911.03840;%%
  %20 citations counted in INSPIRE as of 08 Oct 2020


%\cite{Vladimirov:2020ofp}
\bibitem{Vladimirov:2020ofp} 
  A.~A.~Vladimirov and A.~Schäfer,
  %``Transverse momentum dependent factorization for lattice observables,''
  Phys.\ Rev.\ D {\bf 101}, no. 7, 074517 (2020)
  doi:10.1103/PhysRevD.101.074517
  [arXiv:2002.07527 [hep-ph]].
  %%CITATION = doi:10.1103/PhysRevD.101.074517;%%
  %15 citations counted in INSPIRE as of 08 Oct 2020


%\cite{Ebert:2019tvc}
\bibitem{Ebert:2019tvc} 
  M.~A.~Ebert, I.~W.~Stewart and Y.~Zhao,
  %``Renormalization and Matching for the Collins-Soper Kernel from Lattice QCD,''
  JHEP {\bf 2003}, 099 (2020)
  doi:10.1007/JHEP03(2020)099
  [arXiv:1910.08569 [hep-ph]].
  %%CITATION = doi:10.1007/JHEP03(2020)099;%%
  %23 citations counted in INSPIRE as of 08 Oct 2020


%\cite{Shanahan:2020zxr}
\bibitem{Shanahan:2020zxr} 
  P.~Shanahan, M.~Wagman and Y.~Zhao,
  %``Collins-Soper kernel for TMD evolution from lattice QCD,''
  Phys.\ Rev.\ D {\bf 102}, no. 1, 014511 (2020)
  doi:10.1103/PhysRevD.102.014511
  [arXiv:2003.06063 [hep-lat]].
  %%CITATION = doi:10.1103/PhysRevD.102.014511;%%
  %15 citations counted in INSPIRE as of 08 Oct 2020


%\cite{Bruno:2014jqa}
\bibitem{Bruno:2014jqa} 
  M.~Bruno {\it et al.},
  %``Simulation of QCD with N$_{f} =$ 2 $+$ 1 flavors of non-perturbatively improved Wilson fermions,''
  JHEP {\bf 1502}, 043 (2015)
  doi:10.1007/JHEP02(2015)043
  [arXiv:1411.3982 [hep-lat]].
  %%CITATION = doi:10.1007/JHEP02(2015)043;%%
  %163 citations counted in INSPIRE as of 08 Oct 2020


%\cite{Shanahan:2019zcq}
\bibitem{Shanahan:2019zcq} 
  P.~Shanahan, M.~L.~Wagman and Y.~Zhao,
  %``Nonperturbative renormalization of staple-shaped Wilson line operators in lattice QCD,''
  Phys.\ Rev.\ D {\bf 101}, no. 7, 074505 (2020)
  doi:10.1103/PhysRevD.101.074505
  [arXiv:1911.00800 [hep-lat]].
  %%CITATION = doi:10.1103/PhysRevD.101.074505;%%
  %16 citations counted in INSPIRE as of 08 Oct 2020


\bibitem{supplemental}
Supplemental materials.

%\cite{Edwards:2004sx}
\bibitem{Edwards:2004sx} 
  R.~G.~Edwards {\it et al.} [SciDAC and LHPC and UKQCD Collaborations],
  %``The Chroma software system for lattice QCD,''
  Nucl.\ Phys.\ Proc.\ Suppl.\  {\bf 140}, 832 (2005)
  doi:10.1016/j.nuclphysbps.2004.11.254
  [hep-lat/0409003].
  %%CITATION = doi:10.1016/j.nuclphysbps.2004.11.254;%%
  %738 citations counted in INSPIRE as of 08 Oct 2020

\end{thebibliography}
